% ============================================================================================================
% Optional additions to Section 5 (Results). They add figures, tables and one robustness subsection; they do not
% reword the existing text. Insert each block where its comment says. Upload report/figures/paper_auc_part1.png and
% report/tables/*.tex beside Written_Report.tex (made by scripts/report_figures.py and scripts/report_tables.py).
% Every number below was checked against results/ (see report/factcheck_draft.md).
% ============================================================================================================

% ---- (1) After the paragraph that ends "...than was observable from the participants' characteristics before meeting."
Figure~\ref{fig:auc} and Table~\ref{tab:models} show the held-out AUC of each model with and
without the in-date ratings. Before the date the models are tightly bunched between 0.54 and 0.58; once the ratings
are added they reach 0.78 to 0.85.

\begin{figure}[h]
\centering
\includegraphics[width=\textwidth]{paper_auc_part1.png}
\caption{Before the date every model is near chance; the four-minute date adds most of the signal. Held-out AUC of the
models of Section~4 under event 5-fold cross-validation (4,184 dyads), with 90\% intervals (AUC $\pm\,t_{20}$ times the
event-clustered standard error). Source: results/part1\_models\_metrics.csv.}
\label{fig:auc}
\end{figure}

\input{models}

% ---- (2) After the paragraph that ends "...beyond the participants' individual characteristics."
Table~\ref{tab:ladder} lists each step.

\input{ladder}

% ---- (3) After the paragraph on attractiveness (end of the current Results section).
Table~\ref{tab:stated} gives the stated and revealed importance of all six traits.

\input{stated_revealed}

% ---- (4) A new subsection at the end of Section 5.
\subsection{How robust are these results?}

Two of our conclusions are firm and two are fragile.

\emph{Firm.} Pre-date information carries a small but real signal: every model's 90\% interval in
Figure~\ref{fig:auc} lies above 0.5. And the in-date ratings add far more than any pre-date model achieves: the gain
of about 0.26 in AUC is more than ten times the standard error of any single AUC.

\emph{Fragile.} First, the exact size of the pre-date signal depends on arbitrary choices. Refitting the random
forest with five different random seeds gave pre-date AUCs from 0.5732 to 0.5870, and dealing the events into folds
with ten other seeds moved the logistic model's AUC with a standard deviation of 0.0127. The best pre-date AUC was
0.5796 on complete cases, and 0.560 to 0.579 when the waves with the different preference survey (6--9), or wave 12,
were left out. The pre-date AUC is therefore better described as about 0.56 to 0.59 than as 0.5825 exactly. Second,
the random forest's advantage over logistic regression, 0.0401 under our main specification (90\% interval 0.006 to
0.074, from the event-clustered standard error of the difference), shrank to 0.0215 on complete cases.

Finally, the standard errors themselves were checked by simulation before we used them. With 21 events, our
clustered AUC standard errors were 0.84 to 0.95 of the true spread, so the intervals in Figure~\ref{fig:auc} behave
like roughly 83\% rather than 90\% intervals.
% Sources: results/part1_models_hypotheses.csv, part1_additions_hypotheses.csv, fold_seed_noise_summary.csv, power_summary.csv.
